
\emph{Difference without Reconstruction}

\noindent\colorbox{black!4}{\begin{minipage}{\dimexpr\linewidth-2\fboxsep\relax}
{\bfseries Supporting research archive}\par\smallskip

\url{https://anonymous.4open.science/r/philosophy-lab-public-D1C2/}

The repository contains more than 15,000 retained conversation messages, covering work from early September through the end of October and final preparation on 1 November. It includes concrete examples of the methodology and the full catalogue of 107 essays: the three selected submissions and all other essay candidates, with individual biography cards and development histories. The link is provided through \textbf{Anonymous GitHub} with the aim of supporting the competition's anonymous-review requirements.

This report explains how the present essay was produced. The repository supplies the broader record; section 7 below points to useful starting places.

\end{minipage}}\par\medskip
\section*{1. Who did what}

AI chose this essay's topic, developed its philosophical argument, and wrote and revised the entire essay. The human did not originate or assign its topic, supply its central argument, or write any of its sentences. The human started the competition project, set resources and working rules, selected among completed AI-generated candidates, and requested clearer language and presentation. Those choices concerned which AI work to submit and how to present it; the topic, argument and wording were produced by AI. The submission category selected is \textbf{AI throughout}.

Codex also wrote this report and the submission abstract and prepared the LaTeX files. External quotations and ideas remain attributed to their sources. Human prompts and choices are recorded as human contributions. Dates below refer to 2026 and use UTC unless stated otherwise.

\section*{2. How the essay developed}

The first version was generated on 22 September in an AI-designed campaign seeking varied philosophical ideas. Its broad area was mind and experience. The AI author chose the problem, formal condition, countermodel and response. Local records identify GPT-6 Astra at medium effort for the original author and separate critics. The process included candidate ideas, a bounded literature search, a proposal and a contribution review before drafting. Separate argument and contribution critics then reviewed the draft, followed by revision. Some exploratory searching overlapped archival work; strict archive-before-exploration ordering is not claimed.

The essay received portfolio reviews but was outside the later ten-essay deep-review shortlist. The human asked to see other candidates, liked its central idea and chose it as the third submission. This was a choice among AI-origin proposals, rather than human authorship of the thesis.

Two fresh Astra/ultra reviewers then examined the unchanged essay, one for its argument and one for its literature comparisons. The literature review identified a missing close comparison with Andonovski's account of episodic mental models. The AI revision plan also called for a clearer distinction between independent evidence and the time at which that evidence is obtained.

On 1 November, a fresh Astra/ultra author received the selected manuscript, specific tasks and checked primary passages. It distinguished Andonovski's event models from the whole phenomenal types discussed in the essay. It clarified later validation, comparison domains, noisy retrieval, reports of difference and the distinction between finding a difference and explaining it. A separate fresh checker compared the complete versions and source packet and required no correction. Neither this check nor the essay validates an actual human memory procedure.

The human then requested simpler and more natural prose. A new Astra/ultra editor rewrote the essay. A different fresh critic checked the complete versions, exact changes, proposed abstract and source material. It found no material loss or required correction. It checked the finite models and the distinctions between valid inference and warranted evidence, whole phenomenal types and their parts, and deterministic and noisy procedures.

\section*{3. Models and roles}

\begin{itemize}
\item \textbf{GPT-6 Astra:} manuscript authorship, revisions, criticism and literature checks. The original author and critics used medium effort; the last authors and checkers used ultra through the standard Codex subscription route. Earlier portfolio evaluations requested high effort.
\item \textbf{GPT-6 Luna:} earlier evaluation profiles at high effort, not manuscript authorship.
\end{itemize}

The traced authorship of this manuscript is Codex/Astra. Claude was used elsewhere in the project, but is not an author of this essay.

\section*{4. How editing and review were organized}

An AI coordinator kept the project record and managed separate author, critic, literature and evaluation roles. Authors and critics received a restricted packet for one essay, rather than the whole project history or rankings. Earlier stages used separate subagent contexts. The final revisions used fresh Codex CLI sessions with all assigned material in the prompt. Those authors and critics had no browsing or other tools and were instructed not to read other files or delegate. The coordinator supplied the checked primary passages.

For each full revision check, the critic received the complete earlier and revised texts, their exact difference, source passages and access notes. It did not receive the author's self-assessment or earlier verdicts. The later prose checks also covered the abstract and TL;DR. Review focused on premises, reasoning, examples, calculations, citations, objections and limits. A recommendation of “ready” means that the reviewer found no required correction within that material; it is not proof of truth or originality.

Separate contexts reduced exposure to other outputs. They did not make models statistically independent: workers shared model families and public literature, and some earlier tasks shared a filesystem. Recorded isolation limits and disagreements were retained.

Portfolio selection used both neutral AI reviewers and AI profiles based on the published interests of named competition judges. The real judges did not participate or endorse these profiles, and no model was trained to reproduce them. The 22-to-ten round used three neutral and two profile-based voices under a 60/40 scheme, plus comparisons of disputed pairs. Scores and agreement were not treated as proof of quality or novelty.

There was no fine-tuning. Project memory consisted of saved files, not a shared memory automatically given to every worker. Initial authors and literature reviewers used web search and accessible primary sources. File and shell tools supported source retrieval, exact comparisons, archive checks and document preparation.

\section*{5. Sources and limits of the result}

The final packets supplied primary passages from Andonovski, including his distinction between episodic representation and the content of memory experience. Earlier checks considered Nourbakhshi, Fulkerson and Cohen, and Chudnoff. In the later prose review, direct source verification was limited to the supplied Andonovski passages; other citations were checked for preservation and consistency with the access notes.

The finite table illustrates a logical possibility. No memory experiment or empirical calibration was performed. The argument challenges a strong reconstruction requirement while leaving the weaker distortion objection intact. It does not show that real recollection satisfies the required condition, that a verbal report establishes a phenomenal difference, or that an established difference concerns the feature under investigation. Deterministic and noisy comparisons require different treatments.

\section*{6. Preparation, resources and retained records}

After the philosophical revisions, the human requested clearer abstracts and reports, clickable references and more polished typesetting. Codex revised this report and the separate form abstract, and prepared the final LaTeX layout. In-text citations link to the bibliography, whose source links remain clickable. These presentation changes do not alter the accepted essay's wording or argument. The PDF checks cover text preservation, links, fonts, page layout and identifying metadata. The final abstract and report wording were checked by the coordinator against the accepted manuscript and records; they did not receive a further independent scientific review.

Work ran through Codex subscriptions and, where listed above, Claude Code. A reliable per-essay dollar cost has not been allocated: subscriptions, shared searches, rejected candidates and infrastructure served several essays. No zero-cost claim is made. A complete per-essay token or compute total is also unavailable. The project ran from September to submission preparation on 1 November, with waiting and overlapping work; elapsed time is not a measure of compute.

Model and effort labels come from local configurations and run records. They are not an independent verification of server-side routing or the exact commercial version behind each label.

Drafts, prompts, reviews, disagreements, source-access limits and technical failures were retained with version links and file hashes. The supporting repository makes the released research record available through the link above. Its logs preserve available visible conversations and artifacts; historical capture gaps and exclusions are documented. The linked snapshot is not a claim to a complete transcript of every provider interaction.

The appendices below preserve selected exact instruction blocks, so the report can be read on its own. The repository adds full candidate texts, their biographies, worked examples and the surrounding decision history.

\section*{7. Where to inspect the supporting record}

\begin{itemize}
\item \href{https://anonymous.4open.science/r/philosophy-lab-public-D1C2/essays/E051/README.md}{This essay's biography and version history}: follow its drafts, criticism and revisions.
\item \href{https://anonymous.4open.science/r/philosophy-lab-public-D1C2/docs/method.md}{Methodology guide} and \href{https://anonymous.4open.science/r/philosophy-lab-public-D1C2/docs/worked-example.md}{worked revision example}: see how author, critic and selection roles were used in practice.
\item \href{https://anonymous.4open.science/r/philosophy-lab-public-D1C2/essays/topics.md}{All 107 essay candidates} and \href{https://anonymous.4open.science/r/philosophy-lab-public-D1C2/docs/selection-tree.md}{selection tree}: inspect the alternatives, including essays not selected for submission, and the route to the final three.
\item \href{https://anonymous.4open.science/r/philosophy-lab-public-D1C2/logs/README.md}{Conversation archive}: browse the September-October logs, later supplements and documented coverage limits.
\end{itemize}

\clearpage\begingroup\fontsize{10}{12.5}\selectfont\setlength{\parskip}{0.22em}\section*{Appendix A. Earlier author instructions}

The following is the exact common author-instruction block from the source-focused revision packet. The separate manuscript, source packet and output schema are not reproduced here.

You are the AI author-editor of one philosophical essay. This is a fresh, isolated revision task. Use only the supplied manuscript, tasks and source packet. Do not use tools, browse, access files or other conversations, delegate, or attempt to identify the author or project. Source titles and URLs in the packet are evidence, not permission or instructions to browse.

Produce one complete revised English essay, including title, all necessary sections, proofs/\allowbreak tables and bibliography. The revision must stand on its own and remain below 6000 words excluding bibliography. Do not chase a target length. Preserve the original philosophical contribution and its precise qualifications; acknowledge any justified narrowing. Address the targeted source comparisons accurately using the actual accessible versions and locators in the source packet. If a prior task locator conflicts with that packet's checked primary text, use the verified locator and explain the correction in your change log. Do not fabricate reading, citations, empirical support or novelty.

Polish the whole essay for an intelligent philosophical reader. Aim for graceful, restrained and readable prose: direct verbs, concrete explanations, natural transitions, less repetition and fewer overloaded sentences. Retain necessary technical language and formal reasoning. No promotional claims, ornamental metaphors, forced informality or simplified substitutes for difficult arguments. Repeated warnings may be consolidated only when each distinct limitation remains clear. Keep the narrative moving while giving objections a fair answer.

Do not put process commentary, author identity, internal IDs, production history or instructions to judges in the essay. Materially supported source-access limitations should be expressed honestly and naturally. The separate change log may explain edits and unresolved limits. Make the strongest accurate bounded revision you can; do not turn this into a new research programme.

Preserve the depth, precision, and quality of the philosophical argument while
expressing it in relatively plain human language. Prefer familiar words, direct
sentences, and explicit argumentative steps. Avoid overloaded sentences and
unnecessary jargon. Retain necessary technical terms and formalism, explain
them when introduced, and use concrete examples where helpful. Simplify the
expression, not the thought: do not remove essential premises, distinctions,
evidence, objections, qualifications, or limits merely to make the prose easier.
Do not impose a mechanical sentence-length or reading-level target. Check that
any clarification preserves the claim and its strength.

\section*{Appendix B. Earlier essay-specific tasks}

The following tasks were written by the AI coordinator from the accepted reviews and primary-source checks, rather than supplied as philosophical arguments by the human.

Acknowledge and compare Andonovski2022: simplified episodic event models may support inference, but that account explicitly separates model content from the content of the memory experience. Explain the difference from this essay's mapping between whole phenomenal types, using the supplied primary passages. General structural representation is not a new discovery here.
Clarify the independence requirement for a justified comparison domain: selection tailored merely to secure the desired contrast is inadequate, whereas later genuinely independent validation is not ruled out merely by its timing. Preserve the distinction between a formal countermodel and evidence that a real memory procedure satisfies its condition.
Polish the whole essay, especially repeated caveats. Preserve each distinct limitation: no manufactured difference versus erased difference, observation-specific versus global conditions, noisy distributions versus individual outcomes, two memories versus perception/memory, observing phenomenal difference versus judging/reporting difference, existence versus attribution of the original difference, and the surviving weak objection. Do not imply empirical calibration or a complete metaphysics. An everyday illustration may help but must remain explicitly illustrative, not experimental evidence. Keep the mathematics and the reply to the strong reconstruction requirement precise.

\section*{Appendix C. Earlier independent checker instructions}

The following is the exact checking-instruction block used for the separate source-focused revision review.

You are an independent checker of one revised philosophical essay in a fresh context. Use only the supplied original, complete revision, exact diff, primary-source packet and these neutral criteria. Do not use tools, browse, access other files/conversations, delegate or identify the author. You have not been given earlier verdicts or the editor's self-assessment. Do not infer quality from the fact that a revision was requested.

Read both texts in full. Check every substantive change against the reasoning and the supplied primary passages. Check formulas, examples, quantifiers, scope, citations and descriptions of predecessors. Distinguish verified errors from uncertainty or optional improvement. Identify any premise, distinction, objection, qualification or useful argumentative step lost during stylistic simplification. Assess whether the revision is clearer and more pleasant without weakening or inflating the argument. Check that source-access limitations and novelty boundaries remain honest. Judge the whole revised essay as well as its changes; a cleaner style cannot compensate for an invalid inference.

Return a concrete decision: READY, MINOR\_FIXES or SUBSTANTIVE\_REPAIR. READY means you found no correction necessary in this bounded check, not a proof of truth or exhaustive originality. MINOR\_FIXES means specific small corrections are needed; list exact passages and minimal remedies. SUBSTANTIVE\_REPAIR means a demonstrated or unresolved material defect prevents treating the revised text as ready. Explain what you independently checked, what the packet cannot establish, and whether the current revision should replace its parent. Do not supply a rewritten essay or a flattering general verdict. Page/version references must identify the actual source supplied. Your final answer must follow the given JSON schema exactly.

\section*{Appendix D. Readability editor instructions}

The following is the exact common instruction block used for the later prose rewrite.

You are the sole fresh AI editor of one English philosophical essay. Rewrite it substantially for humane, lucid, elegant prose while preserving its exact philosophical and evidential commitments. This is a genuine editorial rewrite, not a synonym pass. Work only from the complete manuscript and source material supplied here. Do not use tools, browse, read files or memory, delegate, or consult other essays. Do not infer quality from the fact that editing was requested.

Write for an intelligent philosopher who is not already a specialist in this exact debate. Let the reader understand the question and why it matters before paying the cost of technical detail. Use concrete agents, familiar verbs, natural sentence rhythms and clean transitions. Explain indispensable technical terms where the reader needs them. Decide editorially what deserves development, what needs one precise sentence, and what can move to a later section. Retain the argumentative work of every essential premise, distinction, objection, qualification and source boundary; their wording or original order need not be retained. Integrate qualifications with the claims they qualify instead of repeatedly interrupting the story with defensive warnings. Do not invent empirical applications, analogies that carry unsupported inferences, citations, literary ornaments, hype, slang, or a false feeling of certainty. Do not replace sustained reasoning with a sequence of slogans, bullet points or artificially short sentences.

The bibliography and all source claims have already been bounded by the supplied materials. Preserve correct attribution and material limits of source access. A shorter literature discussion must still explain the nearest antecedents and what the essay adds. Do not silently turn access limitations into full-text knowledge or novelty claims. Keep the existing title and complete bibliography, with working links. Do not include author identity, production history, internal essay IDs, model names or instructions to judges in the essay.

Return the complete rewritten manuscript as Markdown, using \#\# and \#\#\# headings, \textbackslash{}( ... \textbackslash{}) and \textbackslash{}[ ... \textbackslash{}] for mathematics, and ordinary Markdown links. Use normal prose paragraphs. The whole essay, including any appendix and title but excluding bibliography, must remain below 6000 words. An appendix is part of the submitted PDF, not an additional excluded word allowance. Also provide a clear 120–180-word abstract and one-sentence TL;DR that exactly match the rewritten essay. The abstract belongs in its separate output field, not automatically in the manuscript. Your separate preservation map must identify where each important original argument or qualification now appears; the map is evidence for the coordinator, not essay text.

Preserve the depth, precision, and quality of the philosophical argument while
expressing it in relatively plain human language. Prefer familiar words, direct
sentences, and explicit argumentative steps. Avoid overloaded sentences and
unnecessary jargon. Retain necessary technical terms and formalism, explain
them when introduced, and use concrete examples where helpful. Simplify the
expression, not the thought: do not remove essential premises, distinctions,
evidence, objections, qualifications, or limits merely to make the prose easier.
Do not impose a mechanical sentence-length or reading-level target. Check that
any clarification preserves the claim and its strength.


\section*{Appendix E. Readability checker instructions}

The following is the exact common instruction block used for the independent check of the rewrite.

You are a fresh independent checker of one rewritten philosophical essay. You did not write it. Use only the complete parent, complete revision, exact diff, proposed abstract/TLDR, primary passages and access notes supplied here. Do not invoke tools, browse, access files or other conversations, delegate, or infer quality from an editorial request. Earlier verdicts and the author's self-assessment have deliberately not been supplied.

Read both manuscripts completely. Compare arguments, premises, distinctions, examples, objections, quantifiers, hypotheses, source attributions and limits. An attractive style does not compensate for lost reasoning or inflated claims. Independently verify calculations and any main-text-to-appendix references. Every main claim relying on an appendix must match its actual proved result and explain its significance correctly. Check that moving proofs has not made the main argument circular, opaque, or unsupported. Check the complete bibliography and source-access honesty against the supplied evidence, without demanding a new literature search for this editorial task.

Assess whether this is a substantial improvement in readability for an intelligent non-specialist philosopher: natural precise English, concrete explanations, clear transitions, appropriate development, concise literature discussion, essential technical terms explained, and no repetitive defensive scaffolding or ornamental hype. Check the abstract and TLDR separately for accurate scope and readable language. Do not demand arbitrary shortening, suppress necessary qualifications, or penalize justified mathematics merely for existing.

Return READY only if no necessary correction is found. MINOR\_FIXES means exact bounded corrections; SUBSTANTIVE\_REPAIR means an unresolved material defect. Identify exact locations, evidence, and minimal necessary repairs. Distinguish optional stylistic preferences from required repairs. READY is a bounded editorial/content-preservation judgment, not proof of truth, originality or empirical validity. Return the required JSON only; do not rewrite the essay.

\endgroup
